\notepage{44}{Rotational invariance and compatible observables}
Rotational invariance can be expressed through commutators:
\begin{align*}
 U_{\hat{\mathbf n}}^\dagger HU_{\hat{\mathbf n}}=H\ \text{for all }\hat{\mathbf n}
 &\iff\comm{H}{U_{\hat{\mathbf n}}}=0\ \text{for all }\hat{\mathbf n}\\
 &\iff\comm{H}{L_i}=0\quad(i=x,y,z).
\end{align*}
For an infinitesimal coordinate-axis rotation this follows from
\[
 0=\comm{H}{I-\ii\epsilon_iL_i/\hbar}
 =-\frac{\ii\epsilon_i}{\hbar}\comm{H}{L_i}.
\]
Writing $L^2=L_x^2+L_y^2+L_z^2$, it also follows that
$\comm{H}{L^2}=0$.

\smallheading{Example: the free particle in three dimensions}
\[
 H=\frac{\mathbf P^2}{2m},\qquad
 E=\frac{p_x^2+p_y^2+p_z^2}{2m}.
\]
Many momentum states have the same energy. Since
\[
 \comm{H}{P_x}=\comm{H}{P_y}=\comm{P_x}{P_y}=0,
\]
one may choose simultaneous eigenstates of $H,P_x,P_y$.
\editorial{These three labels alone do not always remove degeneracy:
at fixed $E,p_x,p_y$, both signs of
$p_z=\pm\sqrt{2mE-p_x^2-p_y^2}$ are possible when the square root is
nonzero. The full momentum vector provides an unambiguous alternative.}

\smallheading{An angular-momentum choice}
Different components of $\mathbf L$ are incompatible:
$\comm{L_i}{L_j}\ne0$ for distinct components in general. However,
$L^2$ is rotationally invariant, and
\[
 \comm{L^2}{L_i}=0.
\]
Choosing the $z$ axis conventionally,
\[
 \comm{H}{L^2}=\comm{H}{L_z}=\comm{L^2}{L_z}=0.
\]
Thus a simultaneous basis may be labeled $\ket{E,\alpha,\beta}$, where
$\alpha$ and $\beta$ are the eigenvalues of $L^2$ and $L_z$.
Additional degeneracy labels are understood when required.

The same compatible set applies to a rotationally invariant central
Hamiltonian
\[
 H=\frac{\mathbf P^2}{2m}+V(R),\qquad R=\abs{\mathbf R}.
\]
